% PNP REVIEW REPORT::REDTEAM::RT-0003::2026-10-11T04:42:19Z
\documentclass[11pt]{article}
\usepackage[margin=0.86in]{geometry}
\usepackage{amsmath,amssymb,amsthm,booktabs,longtable,hyperref}
\newtheorem{theorem}{Theorem}
\newtheorem{lemma}[theorem]{Lemma}
\newcommand{\SAT}{\mathsf{SAT}}
\newcommand{\Ftwo}{\mathbb F_2}
\title{REDTEAM RT-0003:\\Full Compatibility Rank Does Not Give an Exponential Circuit or Transistor Lower Bound}
\author{P versus NP project review}
\date{11 October 2026}
\begin{document}
\maketitle

\paragraph{Result and strict status.}
The newest COST\_CERTIFICATES report correctly proves that four easy
compatibility families can have full $2^n$ rank while admitting direct
linear-time decision.  This audit strengthens that warning in the circuit
model.  In a fixed bounded-fan-in basis, equality, Horn subset, affine
complement, and a treewidth-one independent product all have compatibility
matrices of exact rank $2^n$ over $\Ftwo$, yet explicit circuits use only
$O(n)$ gates.  A declared static-CMOS macro realization therefore uses only
$O(n)$ transistors.  Thus any generic implication charging rank itself as
gate or transistor count is \textbf{FALSIFIED/RETIRED}.  The weaker
$\log_2\operatorname{rank}=n$ scale is consistent with these upper bounds
and still needs a model-matched transfer theorem before it applies to an
arbitrary SAT circuit.

There is a separate quantifier defect in the phrase ``minimum circuit for a
specific SAT instance.''  For each fixed formula, satisfiability is one
constant bit, so a zero-input circuit needs at most one constant-output gate.
The nontrivial target is one circuit $C_N$ deciding every $N$-bit SAT
encoding, with the basis, fan-in, uniformity and invalid-encoding convention
declared.  A superpolynomial lower bound for such general nonuniform
circuits would prove $\SAT\notin\mathsf{P/poly}$ and hence more than
$P\ne NP$.  No such bound is established here.  P versus NP remains
\textbf{OPEN}.  The proofs below are human-readable, not proof-assistant
verification.

\section{Predecessor publication and exact versions}
The audit started from repository head
\texttt{13240ef2b0043a00fc0bda2d144356452b5505dd}.  The immediate predecessor
artifact
\path{research/review/reports/cost_certificates/2026-10-11_0424_CC-0005.tex}
has blob \texttt{e2241c6fffff1af4f94143d4df28dc471803410b}, is substantive, and contains
both the header
\texttt{PNP REVIEW REPORT::COST\_CERTIFICATES::CC-0005} and the required
final sentinel.  Its code, raw result, log and cost-certificate spine were
also present at the recorded verified blobs.  No recovery request was
needed.

The material peer versions additionally included LB-0004 report blob
\texttt{258fe14120d649ef631fca91a7a9f1feceae166a} as read through the current path
\path{research/review/reports/lower_bound/2026-10-10_0410_LB-0004.tex},
RT-0002 blob \texttt{e637c1a2...}, and the current MODEL\_BRIDGE and
PROOF\_MAP manuscripts.  Ellipses here are display abbreviations, not
identifiers used for a write.

\section{Claim card and exact quantifiers}
The attacked generic implication is
\[
 \forall (R_n)_{n\ge1}\quad
 \operatorname{rank}_{\Ftwo}M(R_n)=2^n
 \Longrightarrow
 \operatorname{CircuitSize}(R_n)=2^{\Omega(n)},             \tag{RANK-CIRCUIT}
\]
or the same statement with physical transistor count in place of Boolean
gate count.  Here $R_n:\{0,1\}^n\times\{0,1\}^n\to\{0,1\}$, $M(R_n)$ is
its $2^n$ by $2^n$ communication matrix, and circuits use the basis
$\{\mathrm{AND}_2,\mathrm{OR}_2,\mathrm{NOT}\}$ with unbounded fan-out.
The complete encoded input length is $N=2n$ bits.  The SAME\_REALIZATION
object is one circuit computing the same relation whose matrix rank is
measured.

The weakest useful positive repair would not be (RANK-CIRCUIT).  It would
specify a restricted circuit or protocol model, prove that every admissible
SAT implementation induces the matrix partition used by the rank argument,
and charge the simulation.  For unrestricted SAT circuits the missing
quantifier is
\[
 \exists \{F_N\}\ \forall \{C_N\}\ \exists^\infty N:
 C_N\text{ decides }\SAT_N\Longrightarrow |C_N|>N^k
\]
for every fixed $k$, or an equivalent superpolynomial formulation.  None of
the current compatibility theorems supplies this all-circuit statement.

\section{Exact counterfamilies}
For $x,y\in\{0,1\}^n$, define
\[
 E_n(x,y)=[x=y],\qquad
 H_n(x,y)=[x_i\le y_i\text{ for all }i],
\]
\[
 A_n(x,y)=[y=x\oplus 1^n],\qquad
 T_{n,k}(x,y)=\prod_{j=1}^{k}[x^{(j)}=y^{(j)}],
\]
where the blocks in $T_{n,k}$ partition the $n$ coordinates.
$E_n$ is expressible by independent 2-CNF equality pairs; $H_n$ is the
Horn implication/subset relation; $A_n$ is an affine diagonal solved by
XOR elimination; and $T_{n,k}$ has a matching primal graph and treewidth
one.

\begin{theorem}[FULL-RANK-LINEAR-CIRCUIT1, PROVED/EXACT]
For every $n\ge1$, each of $E_n,H_n,A_n,T_{n,k}$ has
\[
 \operatorname{rank}_{\Ftwo}M=2^n.
\]
In the declared basis, the following upper bounds hold:
\[
 |E_n|=|T_{n,k}|\le5n,\qquad |H_n|\le3n,
 \qquad |A_n|\le6n.
\]
Using the declared macro costs
\[
 \mathrm{NOT}=2,qquad \mathrm{AND}_2=6,qquad
 \mathrm{OR}_2=6
\]
transistors gives implementation upper bounds
\[
 26n-4,qquad 14n-4,qquad 28n-4,qquad 26n-4
\]
respectively.  These transistor values are implementation upper bounds in
one library, not claims of physical minimality.
\end{theorem}

\begin{proof}
$M(E_n)$ is the identity matrix.  $M(A_n)$ is a permutation matrix.  If
subsets are ordered by their binary masks, $M(H_n)$ has entry one exactly
when the row subset is contained in the column subset.  It is upper
triangular with unit diagonal, hence nonsingular over every field.
Finally $M(T_{n,k})$ is a Kronecker product of identity matrices, so its rank
is the product of the block ranks, namely $2^n$.

For one equality coordinate compute XOR as
\[
 (a\lor b)\land\neg(a\land b),
\]
using four gates.  OR the $n$ mismatch bits with $n-1$ gates and negate once,
for $5n$ gates.  Horn subset uses $n$ negations, $n$ conjunctions, an
$(n-1)$-gate OR tree and one final negation, for $3n$ gates.  Affine
complement adds one negation to each equality-coordinate mismatch, for
$6n$ gates.  The independent product is the same coordinatewise equality
circuit with block labels erased.  Summing the declared macro costs yields
the transistor bounds.
\end{proof}

The theorem is an upper-bound counterexample.  It does not claim that
communication rank is useless in every restricted model.  It proves that
the numerical rank $2^n$ cannot be interpreted as $2^n$ mandatory gates or
transistors.  In particular, the usual logarithmic rank scale is only $n$,
which is compatible with the explicit linear circuits.

\section{Fixed-instance hardwiring and the correct circuit target}
\begin{lemma}[INSTANCE-HARDWIRE1, PROVED/EXACT]
For every fixed encoded formula $F$, the Boolean function with no inputs and
value $\SAT(F)$ has a circuit of size zero when constant rails are allowed,
and size at most one when a constant-output gate is counted.
\end{lemma}
\begin{proof}
Choose the constant $1$ circuit if $F$ is satisfiable and the constant $0$
circuit otherwise.
\end{proof}

This is existential nonuniform hardwiring.  It does not give an efficient
procedure for choosing the constant and therefore does not decide SAT.
However, it makes a per-instance minimum-circuit lower bound vacuous.  The
meaningful Boolean function is $\SAT_N:\{0,1\}^N\to\{0,1\}$ on all
canonical $N$-bit encodings, with invalid strings assigned a fixed value.
The relevant distinctions are:
\begin{description}
\item[] \textbf{One fixed instance:} constant output, no uniform acquisition
claim.
\item[] \textbf{One circuit per input length:} nonuniform circuit complexity,
the $\mathsf{P/poly}$ target.
\item[] \textbf{Uniform circuit generator:} an additional algorithmic
uniformity condition, not implied by small nonuniform circuits alone.
\item[] \textbf{Physical chip:} additionally requires a technology, layout,
fan-out, wire, clock, area, delay and power model.
\end{description}

Under $P=NP$, a deterministic polynomial-time SAT algorithm exists and the
standard time-to-circuit simulation yields a polynomial-size uniform circuit
family.  Therefore a purported unconditional superpolynomial general-circuit
lower bound cannot assume $P=NP$ and then derive a contradiction from an
unproved necessity of compatibility extraction.  The lower bound itself is
the missing separation-strength premise.

\section{Wrappers, encodings and one-bit output}
The minimum circuit ignores optional work.  DEAD adds disconnected gates;
COPY duplicates a subcircuit; DELAY inserts two negations; RECOMP evaluates
the same predicate twice and ANDs the answers; INTERP changes the gate-level
representation.  All preserve the one-bit function and can increase a
particular netlist without increasing minimum circuit size.  Padding inputs
may be left unwired.  Duplicate clauses and alternative 2-CNF, Horn or XOR
encodings can enlarge the source description while the semantic circuit
remains linear in the live coordinates.

Consequently, an execution or netlist property must be covariant under
equivalent implementations before it can be a minimum-circuit resource.
Requiring a circuit to output a compatibility table, proof, witness or state
exposure changes the function from one-bit decision and does not lower-bound
the original circuit without an exact reduction back to it.

\section{Bounded reproducible test}
The command
\begin{verbatim}
timeout 30s python3 experiments/review/redteam_rank_circuit_transfer_test.py
  --max-exhaustive-n 8 --max-rank-n 9
  --output results/raw/review/2026-10-11_redteam_rank_circuit_transfer.json
\end{verbatim}
used a 25-second cooperative alarm and completed in 7.530 seconds.  It
performed 349,520 semantic netlist checks, 1,747,600 wrapper-equivalence
checks, 174,760 ignored-padding checks, and 36 exact $\Ftwo$ rank checks.
At $n=9$, all four matrices had rank 512; the gate counts were 45, 27, 54
and 45, with declared macro transistor upper bounds 230, 122, 248 and 230.
All assertions passed.  The experiment is \textbf{COMPUTATIONAL} finite
validation only; the displayed theorem supplies the asymptotic proof.

The script records no randomized samples.  It evaluates every pair through
$n=8$ and computes exact ranks through $n=9$.  Larger experiments were not
needed and no asymptotic extrapolation from timing was made.

\section{Primary literature and barrier audit}
Shannon's switching-circuit synthesis work studies realizing specified
Boolean behavior by switching networks \cite{shannon}.  Its general counting
perspective concerns the prevalence of hard functions, not a lower bound for
the explicit SAT function.  The transfer ``most functions are hard, hence
SAT is hard'' is a \textbf{FAILED TRANSFER}.

Razborov and Rudich formalize a conditional barrier for large, constructive,
useful properties proving strong general-circuit lower bounds
\cite{razborovrudich}.  This audit proposes no such positive property; its
linear circuits are unconditional upper bounds.  Any future use of
compatibility rank as a broad constructive property must nevertheless state
largeness and usefulness rather than merely cite the barrier slogan.

Williams proves superpolynomial lower bounds for nonuniform $\mathrm{ACC}$,
a restricted constant-depth circuit class \cite{williams}.  That theorem is
a genuine lower bound in its declared model, but transferring it to
unrestricted SAT circuits requires an absent compiler forcing arbitrary
circuits into $\mathrm{ACC}$.  This is a \textbf{FAILED TRANSFER} as a free
step and an \textbf{ANALOGY/CANDIDATE} only if such a compiler is separately
specified.

Baker--Gill--Solovay construct relativized worlds with both equality and
separation of P and NP \cite{bgs}.  The counterfamilies and hardwiring lemma
relativize, and an extensional SAT oracle makes the one-bit decision a single
query.  Nothing here supplies a nonrelativizing hinge.  Low-degree or
algebraic extension access likewise does not invalidate the explicit easy
circuits, so the audit does not cross the algebrization barrier.  The PQ
Schur-square and quadratic-relation prompts are \textbf{FAILED TRANSFERS}
for this target unless an exact same-realization theorem converts their
code invariants into general-circuit size; the easy full-rank controls show
that raw algebraic abundance can coexist with simple decision circuitry.

The electronics correspondence is exact only in the easy direction used
here: each declared Boolean macro has a constant transistor implementation,
so an $O(n)$ gate netlist gives an $O(n)$ implementation upper bound.  A
physical VLSI area-time, routing or fan-out lower bound is model-specific and
does not become an unrestricted circuit or Turing-machine lower bound without
a layout-preserving simulation theorem.  That converse transfer remains
\textbf{OPEN}.

\section{Claim-by-claim verdicts}
\begin{longtable}{p{0.26\textwidth}p{0.18\textwidth}p{0.47\textwidth}}
\toprule Claim/version&Verdict&Reason and smallest repair\\\midrule\endhead
CC-0005 full-rank easy families
&PROVED/EXACT accepted&The rank identities and direct semantic algorithms
are correct.  This audit adds explicit bounded-fan-in circuits and transistor
macro upper bounds.\\
Rank magnitude $2^n$ implies $2^{\Omega(n)}$ gates or transistors
&FALSIFIED/RETIRED&Equality already has identity matrix and a $5n$-gate
circuit.  Restrict the model and prove a same-realization simulation before
using rank.\\
Every exact SAT circuit must expose a compatibility table or quotient
&OPEN / unsupported&One-bit circuits for the easy controls do not expose the
table.  A SAT-specific white-box execution relation is required.\\
Minimum hardware for a fixed SAT instance is hard
&FALSIFIED as stated&The fixed answer is a constant.  Replace the target by
the length-indexed function $\SAT_N$ and declare uniformity.\\
General superpolynomial SAT circuit lower bound
&OPEN&It would imply $\SAT\notin\mathsf{P/poly}$, stronger than merely
$P\ne NP$.  No current project theorem establishes it.\\
LB-0004 fiber-type quotient
&PROVED/EXACT scoped accepted&It bounds an endpoint union once types and
binding are supplied; it does not force circuit materialization of those
types.\\
CC-0005 $\mathsf{WExt}$ next target
&OPEN&It must bind to the actual execution, survive equivalent wrappers, and
avoid requiring more than one-bit decision correctness.\\
\bottomrule
\end{longtable}

\section{CORPUS REVIEWED}
The recursive GitHub tree at the starting head, root main, AGENTS, CODEX,
the repository research skill and both red-team protocol references, all
current task specifications, review protocol and STATUS, all current review
spines, main4/main5, newest builder reports, and current coordination ledgers
were surveyed.  Underlying reports reviewed for this target included
MODELBRIDGE-0002, LB-0004, CC-0005 and RT-0002.

Five additional underlying PROOF\_MAP artifacts were read rather than used
through summaries: SEM-QUOTIENT-ACQUIRE1, RAM-RESIDUAL-AUDIT1,
HEC-RAM-BRIDGE-AUDIT, DAG-SHARE-RECOVERY, and
WEIGHTED-POTENTIAL-COUNTEREXAMPLE.  They independently preserve the same
guard: semantic multiplicity, histories, no-reread state, DAG nodes and
weighted syntactic mass do not become arbitrary-algorithm resources without
a charged transfer.  Repository search covered compatibility, rank,
circuit, transistor, semantic quotient, hardwiring, P/poly, padding and
wrapper variants.

\section{DUPLICATION / CONTRADICTION AUDIT}
CC-0005 already proved the central full-rank versus direct-decision
separation.  The new contribution is not another rank computation: it is the
explicit unrestricted-circuit countermodel, the fixed-instance quantifier
repair, and the declared gate-to-transistor upper-bound ledger.  Earlier
RAM-residual and DAG-sharing audits already rejected semantic multiplicity
as mandatory computation, so this report adopts rather than renames that
principle.

No contradiction with LB-0004 was found.  Full-rank relations can still
have many distinct semantic rows, while LB-0004 remains a conditional union
bound after a verifier exposes fiber types.  Neither result forces an
arbitrary circuit to enumerate rows or fibers.  No contradiction with the
root OPEN status or the retired universal cosine endpoint was found.

\section{WHAT CHANGED RELATIVE TO ROOT MAIN}
Root main already warns that structural rank does not itself create a
SAT-semantic return and leaves arbitrary-algorithm acquisition OPEN.  This
audit makes the circuit interpretation explicit: even full rank $2^n$ is
compatible with $O(n)$ unrestricted bounded-fan-in circuit size, and a
specific-instance circuit is constant by hardwiring.  Root main was not
edited because REDTEAM has no LSPV ownership region.

\section{ADOPTED / REJECTED / RETIRED PRIOR COMPONENTS}
\begin{description}
\item[] \textbf{ADOPTED:} CC-0005's four easy full-rank families and
extensional-extractor warning; LB-0004's scoped fiber quotient; SAME\_REALIZATION;
the one-bit output guard; explicit source, encoding, wrapper and acquisition
charges.
\item[] \textbf{REJECTED:} interpreting a communication-matrix dimension as
a gate count; physical transistor language without a declared library and
layout model; per-instance circuit minimum as the SAT-family target.
\item[] \textbf{FALSIFIED/RETIRED:} the generic transfer from full
compatibility rank to exponential unrestricted circuit or transistor size,
and any claim that optional extraction circuitry is unavoidable merely
because the semantic matrix exists.
\end{description}

\section{STALE OR OVERSTATED CLAIMS FOUND}
The current formal review documents are conservatively OPEN and do not state
the falsified exponential rank-to-gate implication.  The phrase ``minimum
number of transistors for a specific SAT problem'' is nevertheless ambiguous
between a fixed instance, an input-length family, and a physical universal
chip.  Future reports should reserve circuit lower bounds for the
length-indexed SAT function and should report transistor counts only as
technology-specific upper or lower bounds.

\paragraph{Single next target.}
Formalize $\mathsf{WExt}(A,F,Q,\tau)$ only after fixing its circuit
counterpart: does a verifier-observable predicate of a one-bit circuit remain
invariant under DEAD/COPY/DELAY/INTERP/RECOMP and under replacing a direct
affine or 2-SAT circuit by an extensionally equal circuit?  If the admissible
$Q$ sets can be disjoint, retire universal extraction.  If not, prove the
SAT-specific source theorem that forces their intersection and charge its
verification without enumerating the compatibility matrix.  This bridge,
the nonrelativizing and nonalgebrizing hinge, general SAT circuit lower
bounds, and P versus NP remain \textbf{OPEN}.

\begin{thebibliography}{9}
\bibitem{shannon} C. E. Shannon,
\emph{The Synthesis of Two-Terminal Switching Circuits}, Bell System
Technical Journal 28 (1949), 59--98.
\url{https://archive.org/download/bstj28-1-59/bstj28-1-59.pdf}.
\bibitem{razborovrudich} A. Razborov and S. Rudich,
\emph{Natural Proofs}, Journal of Computer and System Sciences 55 (1997),
24--35. \url{https://mit6875.github.io/PAPERS/natural_proofs.pdf}.
\bibitem{williams} R. Williams,
\emph{Non-Uniform ACC Circuit Lower Bounds}, Journal of the ACM 61(1)
(2014). \url{https://people.csail.mit.edu/rrw/acc-lbs-journal-final.pdf}.
\bibitem{bgs} T. Baker, J. Gill, and R. Solovay,
\emph{Relativizations of the $P=?NP$ Question}, SIAM Journal on Computing
4 (1975), 431--442. \url{https://doi.org/10.1137/0204037}.
\end{thebibliography}

% END_OF_VERIFIED_REPORT::REDTEAM::RT-0003
\end{document}
